\section{Asesmen Akhir: UTS, UAS, dan Proyek Akhir}
\label{sec:asesmen-akhir}

\subsection{Kisi-kisi Ujian Tengah Semester}

Ujian Tengah Semester diselenggarakan pada minggu ke-8 dan mencakup
\CPMK{01} serta \CPMK{02}, yaitu materi minggu 1 sampai 7.

\begin{table}[H]
    \centering
    \caption{Kisi-kisi Ujian Tengah Semester.}
    \label{tab:kisi-uts}
    \small
    \begin{tabular}{@{}p{1.5cm}p{6.2cm}p{3.4cm}c@{}}
        \toprule
        \textbf{Sub-CPMK} & \textbf{Indikator yang diuji} &
        \textbf{Bentuk soal} & \textbf{Bobot} \\
        \midrule
        1.1 & Arsitektur BIU/EU dan siklus instruksi &
        Esai dan diagram & 20\% \\
        \addlinespace
        1.2 & Fungsi register dan prediksi \eng{flag} &
        Esai dan analisis & 20\% \\
        \addlinespace
        1.3 & Perhitungan alamat fisik &
        Soal hitungan & 20\% \\
        \addlinespace
        2.1 & \eng{Encoding} dan pembacaan \eng{opcode} &
        Soal analisis & 15\% \\
        \addlinespace
        2.2 & Penentuan mode pengalamatan dan alamat efektif &
        Soal analisis & 15\% \\
        \addlinespace
        2.3 & Direktif assembler dan alur translasi &
        Esai dan \eng{coding} terbatas & 10\% \\
        \bottomrule
    \end{tabular}
\end{table}

\subsection{Kisi-kisi Ujian Akhir Semester}

Ujian Akhir Semester diselenggarakan pada minggu ke-16 dan mencakup
\CPMK{03} serta \CPMK{04} secara komprehensif, dengan penekanan
pada kemampuan menulis dan memperbaiki program.

\begin{table}[H]
    \centering
    \caption{Kisi-kisi Ujian Akhir Semester.}
    \label{tab:kisi-uas}
    \small
    \begin{tabular}{@{}p{1.5cm}p{6.2cm}p{3.4cm}c@{}}
        \toprule
        \textbf{Sub-CPMK} & \textbf{Indikator yang diuji} &
        \textbf{Bentuk soal} & \textbf{Bobot} \\
        \midrule
        3.1 & Menulis program aritmatika dan transfer data &
        \eng{Coding} & 20\% \\
        \addlinespace
        3.2 & Manipulasi bit dan pemilihan instruksi geser &
        \eng{Coding} dan analisis & 15\% \\
        \addlinespace
        3.3 & Penerjemahan struktur kontrol &
        \eng{Coding} & 20\% \\
        \addlinespace
        3.4 & Prosedur berparameter dan \eng{stack frame} &
        \eng{Coding} dan analisis stack & 15\% \\
        \addlinespace
        2.4 & Penggunaan interupsi untuk masukan/keluaran &
        \eng{Coding} & 10\% \\
        \addlinespace
        4.1, 4.2 & Menemukan bug dari tabel \eng{trace} &
        Analisis kasus & 10\% \\
        \addlinespace
        4.3 & Membaca \eng{disassembly} dan usulan optimasi &
        Analisis & 10\% \\
        \bottomrule
    \end{tabular}
\end{table}

\subsection{Pedoman Proyek Akhir}

Proyek akhir adalah asesmen puncak mata kuliah ini. Mahasiswa memilih
satu dari empat pilihan berikut, mengerjakannya secara individual
atau berpasangan, lalu mempresentasikannya pada minggu ke-15.

\begin{enumerate}
    \item \textbf{Kalkulator dengan validasi masukan.} Menerima dua
          bilangan dan satu operator dari pengguna, memvalidasi
          masukan, menghitung hasil, dan menampilkan keluaran.
          Wajib menangani pembagian dengan nol.

    \item \textbf{Pengolah string.} Mencari, mengurutkan, dan
          mengganti bagian string. Wajib memakai prosedur dengan
          parameter dan manipulasi memori melalui mode pengalamatan
          tidak langsung.

    \item \textbf{Manipulasi bit untuk enkripsi sederhana.}
          Mengenkripsi dan mendekripsi berkas atau string dengan
          operasi XOR dan rotasi. Wajib menjelaskan mengapa
          enkripsi tersebut \emph{tidak} aman untuk keperluan nyata.

    \item \textbf{\eng{Utility} pemeriksa memori dan register.}
          Menampilkan isi area memori dan nilai register dalam
          format heksadesimal dan ASCII, menyerupai perintah
          \file{D} dan \file{R} pada DEBUG.
\end{enumerate}

Setiap proyek wajib memenuhi syarat minimum berikut:

\begin{itemize}
    \item ditulis dalam assembly 8086 sintaks MASM/TASM dan dapat
          dirakit tanpa kesalahan;
    \item memuat minimal dua prosedur dengan \eng{parameter passing}
          melalui stack;
    \item memuat komentar yang menjelaskan setiap bagian penting;
    \item disertai laporan singkat berisi tujuan, rancangan,
          kendala, dan cara pengujian;
    \item disertai bukti pengujian, termasuk setidaknya satu kasus
          masukan yang salah.
\end{itemize}

\subsection{Rubrik Penilaian Komprehensif}

\begin{table}[H]
    \centering
    \caption{Rubrik penilaian proyek akhir berbasis OBE.}
    \label{tab:rubrik-proyek}
    \footnotesize
    \begin{tabular}{@{}p{3.1cm}p{3.5cm}p{3.5cm}p{3.3cm}@{}}
        \toprule
        \textbf{Kriteria} & \textbf{Sangat baik (4)} &
        \textbf{Baik (3)} & \textbf{Cukup (2)} \\
        \midrule
        Kebenaran fungsi (\CPMK{03}) & Semua fitur berjalan benar,
        termasuk penanganan kasus salah & Fitur utama berjalan,
        sebagian kasus salah belum ditangani & Fitur utama berjalan
        dengan bantuan \\
        \addlinespace
        Kualitas kode (\CPMK{03}) & Struktur prosedur jelas,
        komentar lengkap, tanpa kode berulang & Struktur cukup jelas,
        komentar sebagian & Kode panjang tanpa prosedur \\
        \addlinespace
        Penggunaan instruksi (\CPMK{02}) & Instruksi dan mode
        pengalamatan dipilih tepat dan efisien & Sebagian besar
        tepat, ada pilihan kurang efisien & Banyak pilihan instruksi
        tidak tepat \\
        \addlinespace
        Pengujian dan \eng{debugging} (\CPMK{04}) & Semua bug
        ditemukan dengan bukti tabel \eng{trace} & Sebagian bug
        ditemukan dengan bukti memadai & Bug ditemukan dengan
        bantuan dosen \\
        \addlinespace
        Laporan dan presentasi & Laporan lengkap, presentasi
        jelas dan meyakinkan & Laporan lengkap, presentasi cukup &
        Laporan kurang lengkap \\
        \bottomrule
    \end{tabular}
\end{table}

Nilai proyek dihitung dari rata-rata kelima kriteria dengan bobot
yang sama. Proyek yang tidak memenuhi syarat minimum pada
Bagian~\ref{sec:asesmen-akhir} tidak dapat memperoleh nilai lebih
dari kategori ``cukup''.
